% Source: tex/chapters/ch06/04_ソフトウェアエンジニアリング運用の基礎.tex
\subsubsection{効果的なテストとトラブルシューティング}

運用はシステムの安定性に責任を持つ。そのため、ソフトウェアは本番配備前に十分に試験されなければならない。手作業の試験は時間がかかり、誤りや漏れが発生しやすく、大規模化しにくい。したがって、単体試験、統合試験、システム試験、受け入れ試験、回帰試験を可能な限り自動化することが重要である。

本番でエラーが見つかった場合、運用担当技術者とソフトウェア技術者は、診断を実行し、問題と解決策を記録し、優先順位を付け、影響を評価する。報告された問題が実際に再現するかを確認することも重要である。大規模ソフトウェアに対して全面的な再試験を毎回行うことは高価であるため、回帰試験や試験網羅の戦略が必要になる。

本番環境での試験は難しい。停止できない重要システムでは、カナリア試験や暗黙リリースなどを使い、一部の利用者や裏側の経路だけで新機能を観察する方法が使われる。
